\section{Discussion and Limits}
\label{sec:discussion}
\subsection{Choosing execution granularity}
The results suggest a practical procedure. Start with native batch execution
under the required complete-path semantics: if its per-hop peaks fit the
analysis budget and fresh output matters, immediate execution is a strong
choice. If aligned releases consume too much headroom, try staggering or a
native hybrid, and shorten the hybrid's horizon when age matters. Choose
butterfly-level scheduling when its measured complete-path behavior meets a
requirement that the coarser alternatives do not, at an acceptable total cost.

This is an empirical design rule, not a deadline guarantee. The four-module,
one-thread result is the clearest practical case: relative to vDSP batch, the
production path reduces typical block concentration from about two-thirds of
the budget to about 15\%. Yet the native hybrid is competitive in several
controlled shapes, and the multithreaded engine changes the ranking. The
execution contract remains useful across all of these choices, but
performance cannot be inferred from the number of resumable butterflies alone.
Because Rack's CPU meter reflects aggregate cost, a path with smaller bursts
can still appear more expensive to a user.

Publication age is an explicit cost. A 30\,ms hop adds about 1.8 frames of a
60\,Hz display before polling and repaint, although that arithmetic does not
establish a perceptual threshold. For latest-frame visualization, a UI or
worker thread may be preferable; for exact-hop producer history, such a design
must specify its capture, queueing, overflow, and ownership policy. Spectre's
history storage and per-column mailboxes serve a different consumer contract
from Fourier's latest curve. These alternatives deserve measurement at
equivalent boundaries before broader recommendations are made.

\subsection{What the evidence can establish}
Native layouts, actual module and engine processing, and a concurrent consumer
connect the controlled analyzer to its host. Both pilots, however, used one
M1 Pro, one binary, one compiler, and one date. This is an early Apple Silicon
laptop, neither the base M1 nor a lower performance bound for Apple hardware.
Replication on independent days, with independent builds, and on other
architectures remains outstanding.

Careful setup and power and sleep checks coexist with residual OS activity.
Snapshots cannot identify the cause of an individual outlier; longer compute
paths may be more exposed to interruption, but that mechanism was not
isolated. Maxima and session ranges retain those observations. The fixed group
order also confounds workload family with time and possible thermal drift.

The stream study lacks a separate throughput pass, a matched comparison of
paced and continuous execution at identical boundaries, core-placement
calibration, and a thread-policy sweep. Because four-thread engine workers
consume substantial CPU even in an empty graph, changes with thread count
cannot be attributed to parallel analysis alone. The background sweep did not
locate a saturation threshold, and the extreme module settings lack matched
paced native controls. These gaps limit causal and reliability claims, even
though every reported comparison is traceable to retained data. Headless
consumption verifies copying, ownership, and logical age; device underruns,
GUI load, and repaint are outside its scope. Native allocator internals remain
opaque, and the unused module storage retained by the native controls
precludes a ranking by total memory.

\subsection{Scope of the contribution}
The analytical result bounds credits for a complete dependency-ordered
analyzer. The conditional operation-cost bound explains how weights could be
chosen, whereas the measurements test implementations on one host; cache
effects, unequal unit costs, interrupts, and unrelated engine work separate
these two claims. The empirically tuned weights and the unweighted band-cache
rebuild are implementation choices, not a proof of time balance.

Time-distributed transforms and cooperative processing chains precede this
work. The specific contribution is to combine analysis-tail accounting,
exact-hop cadence, input lifetime, cancellation, and output ownership in one
contract, and to evaluate it with comparisons that expose placement, kernel,
and host effects. Native sub-FFT leaves could improve the tradeoff and remain
the most important unmeasured granularity. Inverse synthesis and long-filter
convolution need their own bounded preparation and delivery stages, and this
analysis-only study makes no performance claim for them.
